\section{Problem Formulation and Metrics}
\label{sec:problem}

\subsection{Setup}
Let $\mathcal{D}=\{(x_i,y_i)\}_{i=1}^{N}$ be a training set and let
$\mathcal{P}\subset\mathcal{D}$ be the poisoned subset, $|\mathcal{P}|=N_p$.
A model $f=\mathcal{A}(\mathcal{D})$ is produced by a fixed training algorithm
$\mathcal{A}$; in this paper $\mathcal{A}$ is the schedule of
Table~\ref{tab:protocol}. For a target class $y_t$ and a trigger function $T$,
the attack success rate is
\begin{equation}
\ASR(f)=\frac{1}{|\mathcal{D}_{\mathrm{test}}|}
\sum_{(x,y)\in\mathcal{D}_{\mathrm{test}}}
\mathbb{1}\!\left[f(T(x))=y_t\right],
\label{eq:asr}
\end{equation}
and clean accuracy is the usual top-1 accuracy on unmodified test inputs. A
backdoor is considered present when $\ASR$ is high and clean accuracy is close
to that of a model trained without poison. The gold reference is
$f_{\mathrm{gold}}=\mathcal{A}(\mathcal{D}\setminus\mathcal{P})$, trained on the
same data with the true poison removed. All success statements are relative to
$f_{\mathrm{gold}}$, not to the poisoned model, because the poisoned model's
clean accuracy can be slightly higher or lower by chance and would otherwise
make weak methods look adequate.

\subsection{Detectors and Their Operating Points}
A detector is a map $\mathrm{DET}:\mathcal{D}\to 2^{\mathcal{D}}$ producing a
flagged set $S$. Relative to the ground-truth poison,
\begin{equation}
r=\frac{|S\cap\mathcal{P}|}{|\mathcal{P}|},\qquad
p=\frac{|S\cap\mathcal{P}|}{|S|},
\label{eq:det}
\end{equation}
with $p$ defined as zero when $S=\emptyset$. Recall $r$ is the fraction of
poison the removal method will see; precision $p$ is the fraction of the flagged
set that is genuinely poisoned. Both matter, but they enter the removal problem
asymmetrically: a missed poison sample remains in the training set and can
sustain the backdoor, whereas a falsely flagged clean sample is merely removed
from training. The experiments are designed to expose that asymmetry.

\subsection{Controlled Forget Sets}
Rather than threshold a detector score, we construct the forget set directly.
For requested recall $r\in[0,1]$ and precision $p\in(0,1]$, let
\begin{equation}
\fset(r,p)=A\cup B,\quad
A\subseteq\mathcal{P},\ |A|=\mathrm{round}(N_p r),
\label{eq:forgetA}
\end{equation}
\begin{equation}
B\subseteq\mathcal{D}_{\mathrm{clean}}\setminus
\mathcal{C}_{y_t},\quad
|B|=\mathrm{round}\!\left(\frac{N_p r(1-p)}{p}\right),
\label{eq:forgetB}
\end{equation}
where $\mathcal{C}_{y_t}$ is the set of clean training samples of the target
class. The sizes follow from the definition of precision: if $|A|=N_p r$ and
$|A|/|A\cup B|=p$, then $|\fset|=N_pr/p$ and $|B|=N_pr(1-p)/p$. False positives
are drawn from non-target clean data so that the forget set never removes clean
examples of the class the attacker wants to protect, which would confound
removal with class deletion. Both $A$ and $B$ are sampled without replacement
with a seed derived from $(r,p,\text{seed})$, so the construction is
deterministic and reproducible. Algorithm~\ref{alg:forget} states the procedure.

\begin{algorithm}[t]
\caption{Controlled forget-set construction}
\label{alg:forget}
\begin{algorithmic}[1]
\Require training set $\mathcal{D}$, poison mask, $r$, $p$, seed $s$
\Ensure $\fset(r,p)$ with $|\fset|=N_pr/p$
\State $\mathcal{P}\gets$ poisoned indices; $N_p\gets|\mathcal{P}|$
\State $n_A\gets\mathrm{round}(N_pr)$
\State $n_B\gets\mathrm{round}(N_pr(1-p)/p)$
\State $A\gets$ sample $n_A$ indices from $\mathcal{P}$ with seed $s$
\State $C\gets$ clean, non-target training indices
\State $B\gets$ sample $n_B$ indices from $C$ with seed $s$
\State \Return $A\cup B$
\end{algorithmic}
\end{algorithm}

The construction has three properties that matter for interpretation. First,
$|\fset|$ grows as $p$ falls at fixed $r$: with $N_p=2{,}500$, the cell
$(r,p)=(0.90,1.0)$ removes 2{,}250 samples, while $(0.90,0.7)$ removes 3{,}214.
Precision is therefore also a measure of how much clean data the method must
give up. Second, $r$ and $p$ are requested values; the realised values are stored
with each run and differ only by rounding, because the counts are integers.
Third, the construction is deliberately kinder than a real detector in one
respect: it never places a false positive in the target class and it never
places poison in the retain set when $r=1$. Real detectors do both. The grid is
therefore a best-case controlled manipulation of detector error, not a
simulation of a specific detector. Section~\ref{sec:cliff} anchors it to two real
detectors by running the same removal methods on their raw flag sets.

\subsection{Defence Success and the Tolerance $\rstar$}
A removal run is a success when both halves of the following criterion hold:
\begin{equation}
\ASR(f')\leq 5\%\quad\text{and}\quad
\mathrm{Acc}_{\mathrm{clean}}(f')\geq
\mathrm{Acc}_{\mathrm{clean}}(f_{\mathrm{gold}})-2.
\label{eq:success}
\end{equation}
The first half says the backdoor is no longer usable at a rate that would matter
to an attacker; the second says the cure did not cost more than two points of
utility. The two-point budget is a modelling choice; it is wide enough to absorb
seed noise (the gold model itself varies by about $\pm0.1$ points) and narrow
enough to exclude methods that remove the backdoor by damaging the network. We
report the two quantities separately in every table so that a reader who prefers
a different budget can re-derive the verdicts. For each method we then define
\begin{equation}
\begin{aligned}
\rstar=\min\Bigl\{r:\ &\overline{\ASR}(r,1)\leq 5\%\ \text{and}\\
&\overline{\mathrm{Acc}}_{\mathrm{clean}}(r,1)\geq
\overline{\mathrm{Acc}}_{\mathrm{clean}}(f_{\mathrm{gold}})-2\Bigr\},
\end{aligned}
\label{eq:rstar}
\end{equation}
where the bar is the mean over the three seeds, and write $\rstar=\text{--}$ when
no recall in the grid succeeds. We use the mean because the criterion describes a
deployment, not a single training run; the per-seed counts are reported
separately in Table~\ref{tab:success} so that a reader who requires unanimity can
see exactly where the two readings diverge. The tolerance is a property of a
removal method, not of a detector: it states how much detector recall the method
needs before it can be relied on.

\subsection{Secondary Metrics}
For the amplification study we ask whether partial removal makes missed poison
separable from clean data. Let $\mathcal{M}=\mathcal{P}\setminus\fset$ be the
poison missed by the forget set and let $\mathcal{C}_0$ be clean
training data of the target class. We fit a detector score to distinguish
$\mathcal{M}$ from $\mathcal{C}_0$ and report the area under the ROC curve,
\begin{equation}
\mathrm{AUC}=\Pr\!\left[s(x_{\mathrm{missed}})>s(x_{\mathrm{clean}})\right],
\label{eq:auc}
\end{equation}
with ties counted at one half. An AUC of 0.5 is chance. For durability we report
ASR after $k$ gradient steps on fresh poison, a loss-threshold membership
inference AUC following Shokri \etal{} \cite{mia}, and the Kullback--Leibler
divergence between the output distributions on clean and trigger-stamped inputs.
A live backdoor shows a large clean-versus-trigger divergence; a dormant one does
not.

\subsection{Statistical Treatment}
Every CIFAR-10 cell is run with seeds 0, 1, and 2. Poison placement, weight
initialisation, and data order vary with the seed; the training and test splits
do not. We report mean $\pm$ standard deviation over seeds and, for the success
criterion, the count of seeds that pass. With three seeds the counts are coarse,
and we do not attach confidence intervals to them. When a method is bimodal
across seeds (NegGrad+ is the clearest case), the mean and standard deviation
describe a mixture rather than a stable operating point; we say so in the table
notes instead of hiding it in a single number.
